Vision--language models (VLMs) are commonly evaluated with image--text matching (ITM) benchmarks, where a model has to decide whether an image and a caption match and the incorrect pairs are handcrafted. We call this binary classification of a pair by a model a matching decision. The same decision is made in retrieval, where matching heads, prompted VLMs and multimodal re-rankers re-rank or filter the top candidates returned by an embedding model. These candidates are similar to each other, and telling them apart needs the fine-grained understanding that ITM benchmarks test, so a decision with a higher ITM accuracy should re-rank them better. In this paper, we test this expectation. We use \svo{} as the ITM benchmark and, with its images and captions as candidates, as a retrieval collection. For 100 caption queries and 100 image queries, the top-10 candidates of five embedding models were judged by human annotators, and we re-rank these lists with seven matching decisions. We find that ITM accuracy orders the decisions correctly but says nothing about the size of their effect. The matching heads of BLIP-2 and FLAVA have almost the same accuracy (63.4\% and 62.4\%), yet re-ordering by the BLIP-2 head raises the judged Success@1 by 6 to 8 points while re-ordering by the FLAVA head lowers it by 5 to 7. Every decision loses much of its gain when it is used as a filter at its own threshold, and none of the gains is visible when only the benchmark's annotated pair of a query counts as correct. Our findings suggest that a matching decision meant for retrieval should be selected on judged retrieved pairs and used as a score rather than as a filter.
